% Cascaded Inverter Schematic (2inv.sp)
% Shows all transistors, subcircuit boundaries, measurement sources, and wiring
\documentclass[border=20pt]{standalone}
\usepackage{circuitikz}
\usetikzlibrary{calc,backgrounds}
\newsavebox{\circuitbox}

\begin{document}
\savebox{\circuitbox}{%
\begin{circuitikz}[american, scale=0.8, transform shape, thick]

% ===== Title =====
\node[font=\large\bfseries, anchor=south] at (12,15.5) {Cascaded Inverter Testbench (2inv.sp)};

% ===== vdd and GND Buses =====
\draw[ultra thick, red!60!black] (-1,14.5) -- (23.5,14.5);
\node[above, font=\bfseries\small, red!60!black] at (-0.2,14.5) {VDD};
\draw[ultra thick, blue!50!black] (-1,0) -- (23.5,0);
\node[below, font=\bfseries\small, blue!50!black] at (-0.2,0) {GND};

% ===============================================================
% VDD SUPPLY CHAIN (x=1)
% Path: 0 -> V_DD -> vdd_src -> R_supply -> vdd_meas -> Vmeas_vdd -> vdd bus
% All labels placed on the LEFT to avoid collision with input chain
% ===============================================================
\draw (1,0) to[V, l_=$V_{DD}$] (1,3.5)
            to[R, l_=$R_{supply}$] (1,7)
            to[american voltage source, l_={$V_{m,vdd}$}] (1,10.5)
            -- (1,14.5);
\filldraw (1,3.5) circle(1.5pt); \node[left, font=\tiny] at (0.85,3.5) {vdd\_src};
\filldraw (1,7) circle(1.5pt); \node[left, font=\tiny] at (0.85,7) {vdd\_meas};
\filldraw (1,10.5) circle(1.5pt); \node[left, font=\scriptsize\bfseries] at (0.85,10.5) {vdd};

% ===============================================================
% INPUT CHAIN (x=5)
% Path: 0 -> Vin(PULSE) -> in_src -> R_in -> in_meas -> Vmeas_in -> in
% All labels placed on the RIGHT to avoid collision with VDD chain
% ===============================================================
\draw (5,0) to[V, l=$V_{in}$] (5,3.5)
            to[R, l=$R_{in}$] (5,7)
            to[american voltage source, l={$V_{m,in}$}] (5,10.5);
\filldraw (5,3.5) circle(1.5pt); \node[left, font=\tiny] at (4.85,3.5) {in\_src};
\filldraw (5,7) circle(1.5pt); \node[left, font=\tiny] at (4.85,7) {in\_meas};
\filldraw (5,10.5) circle(1.5pt); \node[above right, font=\scriptsize\bfseries] at (5.05,10.5) {in};

% ===============================================================
% DRIVER INVERTER (INV1, x=9)
% XMP1: PMOS_MEAS subckt — MP1 + Vmeas_p1
% XMN1: NMOS_MEAS subckt — MN1 + Vmeas_n1
% Vmeas_target: target_int -> target
% ===============================================================

% --- MP1 (PMOS, W=WP1, L=L1) ---
\draw (9,14.5) -- (9,13.5);
\draw (9,13.5) node[pmos, anchor=S] (mp1) {};
\node[right, font=\scriptsize\bfseries] at ($(mp1.center)+(0.55,0.15)$) {MP1};
\node[right, font=\tiny, text=gray] at ($(mp1.center)+(0.55,-0.25)$) {W\!=\!WP1, L\!=\!L1};

% Vmeas_p1 (0V source inside PMOS_MEAS)
\filldraw (mp1.D) circle(1.5pt);
\node[right, font=\tiny] at ($(mp1.D)+(0.15,0.15)$) {target\_mp1};
\draw (mp1.D) to[american voltage source, l_={$V_{p1}$}, a^={\tiny 0V}] (9,10);

% Wire to target_int
\draw (9,10) -- (9,7.5) coordinate (target-int);
\filldraw (target-int) circle(2pt);
\node[below left, font=\scriptsize\bfseries] at ($(target-int)+(-0.1,-0.1)$) {target\_int};

% --- MN1 (NMOS, W=WN1, L=L1) ---
\draw (9,0) -- (9,1);
\draw (9,1) node[nmos, anchor=S] (mn1) {};
\node[right, font=\scriptsize\bfseries] at ($(mn1.center)+(0.55,0.15)$) {MN1};
\node[right, font=\tiny, text=gray] at ($(mn1.center)+(0.55,-0.25)$) {W\!=\!WN1, L\!=\!L1};

% Vmeas_n1 (0V source inside NMOS_MEAS)
\filldraw (mn1.D) circle(1.5pt);
\node[right, font=\tiny] at ($(mn1.D)+(0.15,-0.15)$) {target\_mn1};
\draw (mn1.D) to[american voltage source, l={$V_{n1}$}, a={\tiny 0V}] (9,5.5);

% Wire to target_int
\draw (9,5.5) -- (target-int);

% --- INV1 Gate Bus (connects both gates to "in") ---
\draw (mp1.G) -- +(-2,0) coordinate (g1t);
\draw (mn1.G) -- +(-2,0) coordinate (g1b);
\draw (g1t) -- (g1b);
\draw (5,10.5) -- (g1t |- 5,10.5) node[circ]{};

% --- Vmeas_target (top-level: target_int -> target) ---
\draw (target-int) -- (10,7.5);
\draw (10,7.5) to[american voltage source, l_={$V_{m,tgt}$}, a^={\tiny 0V}] (12.5,7.5)
      coordinate (target-node);

% ===============================================================
% LOAD INVERTER (INV2, x=16)
% XMP2: PMOS_MEAS subckt — MP2 + Vmeas_p2
% XMN2: NMOS_MEAS subckt — MN2 + Vmeas_n2
% Vmeas_out: out_meas -> out
% ===============================================================

% --- MP2 (PMOS, W=WP2, L=L2) ---
\draw (16,14.5) -- (16,13.5);
\draw (16,13.5) node[pmos, anchor=S] (mp2) {};
\node[right, font=\scriptsize\bfseries] at ($(mp2.center)+(0.55,0.15)$) {MP2};
\node[right, font=\tiny, text=gray] at ($(mp2.center)+(0.55,-0.25)$) {W\!=\!WP2, L\!=\!L2};

% Vmeas_p2 (0V source inside PMOS_MEAS)
\filldraw (mp2.D) circle(1.5pt);
\node[right, font=\tiny] at ($(mp2.D)+(0.15,0.15)$) {out\_mp2};
\draw (mp2.D) to[american voltage source, l_={$V_{p2}$}, a^={\tiny 0V}] (16,10);

% Wire to out_meas
\draw (16,10) -- (16,7.5) coordinate (out-meas);
\filldraw (out-meas) circle(2pt);
\node[below left, font=\scriptsize\bfseries] at ($(out-meas)+(-0.1,-0.1)$) {out\_meas};

% --- MN2 (NMOS, W=WN2, L=L2) ---
\draw (16,0) -- (16,1);
\draw (16,1) node[nmos, anchor=S] (mn2) {};
\node[right, font=\scriptsize\bfseries] at ($(mn2.center)+(0.55,0.15)$) {MN2};
\node[right, font=\tiny, text=gray] at ($(mn2.center)+(0.55,-0.25)$) {W\!=\!WN2, L\!=\!L2};

% Vmeas_n2 (0V source inside NMOS_MEAS)
\filldraw (mn2.D) circle(1.5pt);
\node[right, font=\tiny] at ($(mn2.D)+(0.15,-0.15)$) {out\_mn2};
\draw (mn2.D) to[american voltage source, l={$V_{n2}$}, a={\tiny 0V}] (16,5.5);

% Wire to out_meas
\draw (16,5.5) -- (out-meas);

% --- INV2 Gate Bus (connects both gates to "target") ---
\draw (mp2.G) -- +(-2,0) coordinate (g2t);
\draw (mn2.G) -- +(-2,0) coordinate (g2b);
\draw (g2t) -- (g2b);
\draw (target-node) -- (g2t |- target-node) node[circ]{};
\node[above left, font=\scriptsize\bfseries] at (g2t |- target-node) {target};

% --- Vmeas_out (top-level: out_meas -> out) ---
\draw (out-meas) -- (17.5,7.5);
\draw (17.5,7.5) to[american voltage source, l_={$V_{m,out}$}, a^={\tiny 0V}] (19.5,7.5)
      coordinate (out-node);
\filldraw (out-node) circle(2pt);
\node[above right, font=\scriptsize\bfseries] at (19.5,7.55) {out};

% ===============================================================
% LOAD CAPACITOR (XCload: C_MEAS subcircuit)
% Path: out -> Vmeas_c -> C_load (5fF) -> GND
% ===============================================================
\draw (out-node) -- (21,7.5) -- (21,6);
\draw (21,6) to[american voltage source, l_={$V_{m,c}$}] (21,2.5);
\filldraw (21,2.5) circle(1.5pt);
\node[right, font=\tiny] at (21.15,2.5) {out\_cload};
\draw (21,2.5) to[C, l_={$C_{load}$}, a^={\tiny 5\,fF}] (21,0);

% ===============================================================
% SUBCIRCUIT BOUNDARIES (dashed boxes)
% ===============================================================
\begin{scope}[on background layer]
  % XMP1: PMOS_MEAS (MP1 + Vmeas_p1)
  \draw[dashed, rounded corners=4pt, gray!50, thick]
    (7.4,9.2) rectangle (11.6,14.2);
  \node[font=\tiny, text=gray!50, anchor=south east] at (11.55,9.25) {XMP1: PMOS\_MEAS};

  % XMN1: NMOS_MEAS (MN1 + Vmeas_n1)
  \draw[dashed, rounded corners=4pt, gray!50, thick]
    (7.4,0.3) rectangle (11.6,5.8);
  \node[font=\tiny, text=gray!50, anchor=north east] at (11.55,5.75) {XMN1: NMOS\_MEAS};

  % XMP2: PMOS_MEAS (MP2 + Vmeas_p2)
  \draw[dashed, rounded corners=4pt, gray!50, thick]
    (14.4,9.2) rectangle (18.6,14.2);
  \node[font=\tiny, text=gray!50, anchor=south east] at (18.55,9.25) {XMP2: PMOS\_MEAS};

  % XMN2: NMOS_MEAS (MN2 + Vmeas_n2)
  \draw[dashed, rounded corners=4pt, gray!50, thick]
    (14.4,0.3) rectangle (18.6,5.8);
  \node[font=\tiny, text=gray!50, anchor=north east] at (18.55,5.75) {XMN2: NMOS\_MEAS};

  % XCload: C_MEAS (Vmeas_c + C_load)
  \draw[dashed, rounded corners=4pt, gray!50, thick]
    (19.2,0.3) rectangle (23.3,5.8);
  \node[font=\tiny, text=gray!50, anchor=north east] at (23.25,5.75) {XCload: C\_MEAS};
\end{scope}

\end{circuitikz}%
}% end savebox
\begin{minipage}[t]{\wd\circuitbox}
\usebox{\circuitbox}

\vspace{1.5em}
\noindent\rule{\linewidth}{0.4pt}
\vspace{1em}


\end{minipage}

\end{document}
